\documentclass[11pt,letterpaper,oneside]{report}
\usepackage[T1]{fontenc}
\usepackage[utf8]{inputenc}
\usepackage{lmodern}
\usepackage[letterpaper,margin=0.92in,headheight=15pt]{geometry}
\usepackage{amsmath,amssymb,amsthm,mathtools}
\usepackage{microtype}
\usepackage[dvipsnames]{xcolor}
\usepackage{graphicx}
\usepackage{booktabs,longtable,array,tabularx}
\usepackage{enumitem}
\usepackage{fancyhdr}
\usepackage{titlesec}
\usepackage[most]{tcolorbox}
\usepackage{listings}
\usepackage{tikz}
\usetikzlibrary{positioning}
\usepackage{hyperref}
\hypersetup{colorlinks=true,linkcolor=MidnightBlue,urlcolor=RoyalBlue,pdfauthor={Independent Mathematics and Engineering Study Suite},pdftitle={Prerequisite-Based Curriculum and Problem Workbook}}
\definecolor{ink}{HTML}{162639}
\definecolor{navy}{HTML}{173B61}
\definecolor{blue}{HTML}{2874A6}
\definecolor{teal}{HTML}{007C83}
\definecolor{amber}{HTML}{8B5E16}
\definecolor{pale}{HTML}{EDF4F8}
\definecolor{paleteal}{HTML}{ECF8F5}
\definecolor{paleamber}{HTML}{FFF6E8}
\titleformat{\chapter}[display]{\normalfont\huge\bfseries\color{navy}}{\filleft\fontsize{43}{46}\selectfont\thechapter}{0.4ex}{\titlerule\vspace{1ex}}[\vspace{1ex}\titlerule]
\titleformat{\section}{\normalfont\Large\bfseries\color{navy}}{\thesection}{0.8em}{}
\titleformat{\subsection}{\normalfont\large\bfseries\color{teal}}{\thesubsection}{0.75em}{}
\pagestyle{fancy}\fancyhf{}\fancyhead[L]{\small\color{navy}Differential Equations Engineering Suite}\fancyhead[R]{\small\color{navy}Volume 8 of 8}\fancyfoot[C]{\thepage}\renewcommand{\headrulewidth}{0.3pt}
\setlist[itemize]{leftmargin=*,itemsep=3pt,topsep=5pt}\setlist[enumerate]{leftmargin=*,itemsep=4pt,topsep=5pt}
\setlength{\parskip}{5pt}\setlength{\parindent}{0pt}\setlength{\emergencystretch}{2em}
\newtheorem{theorem}{Theorem}[chapter]
\newtheorem{proposition}[theorem]{Proposition}
\newtheorem{lemma}[theorem]{Lemma}
\theoremstyle{definition}\newtheorem{definition}[theorem]{Definition}
\newtheorem{example}[theorem]{Worked example}
\theoremstyle{remark}\newtheorem{remark}[theorem]{Remark}
\newtcolorbox{bookground}[1][]{enhanced,breakable,colback=pale,colframe=blue,boxrule=0.6pt,title={TEXTBOOK ALIGNMENT},fonttitle=\sffamily\bfseries\small,#1}
\newtcolorbox{modern}[1][]{enhanced,breakable,colback=paleteal,colframe=teal,boxrule=0.6pt,title={MODERN COMPUTATIONAL EXTENSION},fonttitle=\sffamily\bfseries\small,#1}
\newtcolorbox{caution}[1][]{enhanced,breakable,colback=paleamber,colframe=amber,boxrule=0.6pt,title={ASSUMPTION / VALIDATION},fonttitle=\sffamily\bfseries\small,#1}
\newcommand{\dd}{\,\mathrm{d}}
\newcommand{\R}{\mathbb{R}}
\newcommand{\lap}{\mathcal{L}}
\newcommand{\norm}[1]{\left\lVert#1\right\rVert}
\newcommand{\abs}[1]{\left\lvert#1\right\rvert}
\lstdefinestyle{engpy}{language=Python,basicstyle=\ttfamily\footnotesize,keywordstyle=\color{RoyalBlue}\bfseries,commentstyle=\color{gray},stringstyle=\color{teal},backgroundcolor=\color{pale},frame=single,rulecolor=\color{blue!30},breaklines=true,columns=fullflexible,keepspaces=true,showstringspaces=false,tabsize=4,upquote=true}
\lstset{style=engpy}
\begin{document}
\begin{titlepage}
\pagecolor{navy}\color{white}
\vspace*{0.75in}
{\sffamily\bfseries\large MATHEMATICS \& ENGINEERING DOCUMENTATION SUITE\\[12pt]}
{\color{cyan!40}\rule{\textwidth}{2pt}}\\[1.0in]
{\sffamily\bfseries\fontsize{29}{36}\selectfont Prerequisite-Based Curriculum and Problem Workbook\par}
\vspace{0.45in}
{\Large Volume 8 of 8}\par\vspace{0.55in}
{\color{cyan!30}\rule{0.50\textwidth}{1pt}}\par\vspace{0.3in}
{\large Independent technical companion aligned with the\par}
{\large 13-chapter organization of Nagle, Saff \& Snider (7th ed.)\par}
\vfill
{\large Original derivations, worked examples, verification and engineering practice\par}\vspace{0.3in}
{\small Prepared October 2026 \quad\textbullet\quad pdf\LaTeX{} compatible\par}
\end{titlepage}
\nopagecolor\color{ink}\pagenumbering{roman}
\chapter*{Reader's guide and source boundaries}\addcontentsline{toc}{chapter}{Reader's guide and source boundaries}
This independent study resource is organized around \emph{Fundamentals of Differential Equations and Boundary Value Problems}, Seventh Edition, by R.~Kent Nagle, Edward~B. Saff, and Arthur~David Snider (Pearson, 2018; ISBN 978-0-321-97710-6). The book's chapter and section numbers are given as \emph{navigation references}, not as claims of reproduced text or proofs.

\begin{bookground}
\textbf{Textbook-grounded scope:} topics, terminology, and chapter/section organization corresponding to the cited textbook. All explanations, derivations, exercise prompts, and diagrams in this volume are newly authored and may emphasize different examples. Check the textbook itself for its exact hypotheses, proofs, exercises, or nomenclature in a particular section.
\end{bookground}
\begin{modern}
\textbf{Distinct extension layer:} Python/NumPy/SciPy workflows, engineering QA and reproducibility practices, automated convergence studies, diagnostic plots, and implementation guidance. These are \emph{not} represented as content supplied by the textbook. API calls are examples and should be verified against the installed package version.
\end{modern}
\begin{caution}
\textbf{Mathematical qualification:} formulas require their stated regularity, domain and boundary/initial conditions. Numerical agreement does not replace a proof; computational outputs are not certified error bounds unless explicitly established.
\end{caution}
\tableofcontents\clearpage\pagenumbering{arabic}
\chapter{Learning architecture and prerequisites}
\begin{bookground}The 13 modules below follow the textbook's 13 chapters and preserve their original progression. The schedules, assessments, exercises and coding labs are newly authored instructional guidance; they are not transcriptions of the book's problems or sample syllabi.\end{bookground}
\section{Learning objectives}
At completion, the learner should be able to classify differential equations, solve representative IVPs and BVPs analytically, construct and verify mathematical models, implement and validate numerical methods, analyze linear and nonlinear dynamics, derive separated PDE solutions, use Sturm--Liouville theory, and reason precisely about existence and uniqueness.
\section{Prerequisite DAG}
\begin{center}
\begin{tikzpicture}[node distance=1.2cm, every node/.style={draw,rounded corners,align=center,font=\footnotesize,minimum width=3.3cm,fill=pale},>=stealth]
\node (calc) {Calculus and series};
\node[right=1.1cm of calc] (alg) {Linear algebra};
\node[below=of calc] (first) {Ch. 1--2: First order};
\node[below=of alg] (systems) {Ch. 5, 9: Systems};
\node[below=of first] (second) {Ch. 3--4, 6--8:\\ models, higher order, transforms};
\node[below=of systems] (stability) {Ch. 12: Stability};
\node[below=1.35cm of second] (pde) {Ch. 10--11: PDE / eigenvalue BVP};
\node[below=of pde] (theory) {Ch. 13: Existence / uniqueness};
\draw[->] (calc)--(first);\draw[->](alg)--(systems);\draw[->](first)--(second);\draw[->](first)--(systems);\draw[->](systems)--(stability);\draw[->](second)--(pde);\draw[->](systems)--(pde);\draw[->](first) to[bend right=35] (theory);
\end{tikzpicture}
\end{center}
\begin{caution}The diagram is an instructional dependency recommendation, not a claim that the book mandates this exact ordering. Chapters 5 and 9 can be introduced at different levels of linear algebra preparation.\end{caution}
\section{Workload and evidence of mastery}
A 26-week two-term path allocates approximately two weeks per chapter, plus adaptive review. An accelerated 13-week path can use one module per week but needs substantial prior calculus and linear algebra. Allocate roughly one quarter of study time each to derivations, worked examples, problem-solving, and computational validation. Record evidence through analytic residuals and convergence studies instead of only final answers.
\chapter{Modules 1--3: foundations, first-order methods and models}
\section{Module 1: Introduction (Chapter 1, \S1.1--1.4)}
\textbf{Prerequisites:} differentiation, elementary integration, function graphs. \textbf{Learning targets:} classification, initial values, solution intervals, direction fields, Euler approximation. \textbf{Paper lab:} solve $y'=-2y$, $y(0)=1$, and sketch slopes in the $ty$ plane. \textbf{Code lab:} implement Euler at steps $0.2,0.1,0.05$ and compare to $e^{-2t}$ at $t=1$. \textbf{Mastery evidence:} correctly identify the solution interval and demonstrate first-order global convergence.
\section{Module 2: First-Order Differential Equations (Chapter 2, \S2.1--2.6)}
\textbf{Prerequisites:} Module 1 and integration techniques. \textbf{Learning targets:} separation, linear integrating factors, exact differential equations, special factors and transformations. \textbf{Paper lab:} solve $y'+3y=6$, $y(0)=4$. \textbf{Code lab:} validate solutions by differentiating with SymPy; identify singular denominators. \textbf{Mastery evidence:} method-selection argument, residual zero, excluded branches handled.
\section{Module 3: Mathematical Models and Numerical Methods Involving First-Order Equations (Chapter 3, \S3.1--3.7)}
\textbf{Prerequisites:} Modules 1--2. \textbf{Learning targets:} mass-balance, heat/cooling, mechanics, circuits, Euler and Runge--Kutta. \textbf{Paper lab:} derive the perfectly mixed tank model with constant volume. \textbf{Code lab:} compare Euler, RK4 and a SciPy adaptive solver on identical forcing. \textbf{Mastery evidence:} units, conservation checks, solver settings and step-refinement table.
\chapter{Modules 4--6: linear equations and systems}
\section{Module 4: Linear Second-Order Equations (Chapter 4, \S4.1--4.10)}
\textbf{Prerequisites:} Modules 1--2 and trigonometric/exponential algebra. \textbf{Learning targets:} characteristic roots, complex roots, superposition, undetermined coefficients, variation of parameters and vibration. \textbf{Paper lab:} solve $y''+2y'+5y=0$ with arbitrary initial conditions. \textbf{Code lab:} compare analytic modes to a damped-oscillator numerical integration. \textbf{Mastery evidence:} homogeneous basis and initial conditions verified independently.
\section{Module 5: Introduction to Systems and Phase Plane Analysis (Chapter 5, \S5.1--5.8)}
\textbf{Prerequisites:} first-order IVPs; basic vectors beneficial. \textbf{Learning targets:} compartment systems, numerical state space, phase plane, epidemic and oscillatory models, qualitative behavior. \textbf{Paper lab:} derive conservation for two closed mixing tanks. \textbf{Code lab:} plot phase portrait and assess equilibria. \textbf{Mastery evidence:} clear distinction between a geometric orbit and a time-parametrized solution.
\section{Module 6: Theory of Higher-Order Linear Differential Equations (Chapter 6, \S6.1--6.4)}
\textbf{Prerequisites:} Module 4 and linear independence. \textbf{Learning targets:} fundamental solutions, repeated roots, annihilator method and variation of parameters. \textbf{Paper lab:} solve $(D+2)^3y=0$. \textbf{Code lab:} check the three basis functions with symbolic differentiation. \textbf{Mastery evidence:} correct dimension, basis independence and solution residual.
\chapter{Modules 7--9: transforms, series and matrix methods}
\section{Module 7: Laplace Transforms (Chapter 7, \S7.1--7.10)}
\textbf{Prerequisites:} integrals, exponential functions and linear ODEs. \textbf{Learning targets:} transform definition and properties, inverse transforms, IVPs, discontinuous/periodic forcing, convolution, impulse inputs and systems. \textbf{Paper lab:} solve $y'+y=u(t-1)$ with $y(0)=0$. \textbf{Code lab:} compare direct integration of piecewise forcing with the analytic unit-step solution. \textbf{Mastery evidence:} correct shift argument and causal activation.
\section{Module 8: Series Solutions of Differential Equations (Chapter 8, \S8.1--8.8)}
\textbf{Prerequisites:} Taylor series, convergence basics and second-order ODEs. \textbf{Learning targets:} analytic coefficients, recurrences, Cauchy--Euler, Frobenius, regular singular points and special functions. \textbf{Paper lab:} derive first coefficients of $y''+xy=0$. \textbf{Code lab:} compare truncated-series residuals and a numerical reference. \textbf{Mastery evidence:} two independent free constants and a documented convergence caveat.
\section{Module 9: Matrix Methods for Linear Systems (Chapter 9, \S9.1--9.8)}
\textbf{Prerequisites:} matrices, eigenvalues, Modules 4--5. \textbf{Learning targets:} normal form, constant coefficient systems, complex eigenvalues, nonhomogeneous systems and matrix exponential. \textbf{Paper lab:} compute $e^{At}$ for a diagonal $A$. \textbf{Code lab:} compare \texttt{scipy.linalg.expm} to numerical integration. \textbf{Mastery evidence:} accurate transient modes, attention to defective matrices.
\chapter{Modules 10--13: PDEs, boundary problems, stability and theory}
\section{Module 10: Partial Differential Equations (Chapter 10, \S10.1--10.7)}
\textbf{Prerequisites:} partial derivatives, Fourier series, linear ODE eigenvalue problems. \textbf{Learning targets:} separation, sine/cosine Fourier series, heat, wave and Laplace equations. \textbf{Paper lab:} solve zero-endpoint heat flow with one sine mode. \textbf{Code lab:} simulate a finite-difference heat equation; check stability and energy decay. \textbf{Mastery evidence:} every initial/boundary condition verified.
\section{Module 11: Eigenvalue Problems and Sturm--Liouville Equations (Chapter 11, \S11.1--11.8)}
\textbf{Prerequisites:} integration by parts, eigenvalues, Module 10. \textbf{Learning targets:} regular/singular Sturm--Liouville, orthogonality, Fredholm alternative, eigenfunction expansion, Green's functions, oscillation/comparison. \textbf{Paper lab:} derive eigenvalues for Dirichlet $-y''$. \textbf{Code lab:} compute discrete eigenpairs and a Green-integral benchmark. \textbf{Mastery evidence:} correct weight and compatibility conditions at resonance.
\section{Module 12: Stability of Autonomous Systems (Chapter 12, \S12.1--12.8)}
\textbf{Prerequisites:} Jacobians, eigenvalues and Module 5. \textbf{Learning targets:} planar systems, almost linear systems, energy, Lyapunov direct method, limit cycles, higher-dimensional stability and neuron modeling. \textbf{Paper lab:} choose a Lyapunov function for a damped oscillator. \textbf{Code lab:} plot a nonlinear phase portrait and verify an invariant set. \textbf{Mastery evidence:} valid distinction between stability and asymptotic stability.
\section{Module 13: Existence and Uniqueness Theory (Chapter 13, \S13.1--13.4)}
\textbf{Prerequisites:} analysis of IVPs, continuity, integration, convergence, Lipschitz estimates. \textbf{Learning targets:} successive approximations, Picard's theorem, linear existence, continuous dependence. \textbf{Paper lab:} compute first four Picard iterations for $y'=y$, $y(0)=1$. \textbf{Code lab:} study sensitivity to perturbed initial values and compare with Gronwall bounds. \textbf{Mastery evidence:} hypotheses and conclusions of existence versus uniqueness correctly separated.
\chapter{Textbook chapter and section navigation index}
\begin{bookground}This is an organizational map. It is neither a reproduction of the book's exposition nor a claim that every exercise in the book is solved in this suite.\end{bookground}
\small
\begin{longtable}{p{0.08\textwidth}p{0.39\textwidth}p{0.43\textwidth}}
\toprule Chapter & Book chapter title & Major section progression and linked suite volumes \\ \midrule\endhead
1 & Introduction & Background, solutions/IVPs, direction fields, Euler; Vols. 1, 4 \\
2 & First-Order Differential Equations & Falling body, separable, linear, exact, integrating factors, transformations; Vols. 2, 3 \\
3 & Mathematical Models and Numerical Methods Involving First-Order Equations & Models, compartments, buildings, Newtonian mechanics, circuits, Euler, Runge--Kutta; Vols. 4, 7 \\
4 & Linear Second-Order Equations & Oscillators, homogeneous/forced solutions, variation, variable coefficients, vibration; Vols. 2, 3, 7 \\
5 & Introduction to Systems and Phase Plane Analysis & Interconnected tanks, elimination, numerical systems, phase plane, biomathematics, mechanics/electrical, chaos; Vols. 5, 7 \\
6 & Theory of Higher-Order Linear Differential Equations & Linear theory, constant coefficients, annihilators and variation; Vols. 2, 3 \\
7 & Laplace Transforms & Definition/properties, inverse, IVPs, discontinuous/periodic forcing, convolution, delta, systems; Vol. 3 \\
8 & Series Solutions of Differential Equations & Taylor, analytic coefficients, power series, Euler/Frobenius, special functions; Vols. 2, 3 \\
9 & Matrix Methods for Linear Systems & Matrices/vectors, normal form, eigenvalues, matrix exponential; Vol. 5 \\
10 & Partial Differential Equations & Separation, Fourier sine/cosine series, heat/wave/Laplace; Vols. 6, 7 \\
11 & Eigenvalue Problems and Sturm--Liouville Equations & Regular/singular theory, Fredholm, expansions, Green, comparison; Vol. 6 \\
12 & Stability of Autonomous Systems & Planar systems, energy, Lyapunov, limit cycles, neurons; Vols. 5, 7 \\
13 & Existence and Uniqueness Theory & Picard iteration, existence, continuous dependence; Vol. 1 \\
\bottomrule
\end{longtable}
\normalsize
\chapter{Problem-solving protocol and marking rubric}
\section{Standard written solution protocol}
\begin{enumerate}
\item Restate the equation with all initial/boundary conditions and the domain.
\item Classify the equation, identify a method and justify why it applies.
\item Carry through intermediate algebra and list any division or substitution restrictions.
\item Fit constants and check all conditions. Explicitly verify by substitution.
\item For numerical tasks, specify method, grid/tolerances, residual and uncertainty.
\item Interpret the result in context, including limitations and nonuniqueness.
\end{enumerate}
\section{Assessment rubric}
\begin{longtable}{p{0.43\textwidth}p{0.12\textwidth}p{0.35\textwidth}}
\toprule Criterion & Weight & Evidence \\ \midrule\endhead
Classification and assumptions & 15\% & Domain, order, regularity, constraints \\
Derivation and correct method & 25\% & Reproducible intermediate steps \\
Initial/boundary conditions & 15\% & Independent checks, compatibility \\
Equation residual verification & 15\% & Algebraic or computed residual \\
Numerical convergence and reproducibility & 15\% & Step refinement and settings \\
Interpretation and communication & 15\% & Units, uncertainty, limitations \\
\bottomrule
\end{longtable}
\chapter{Original problem workbook: Part I}
\section{First-order equations and methods}
\begin{enumerate}[label=\textbf{P\arabic*.}]
\item Solve $y'=2t y$ with $y(0)=3$.
\item Solve $y'+2y=6$, $y(0)=0$.
\item Find equilibrium solutions for $y'=y(3-y)$ and classify local directions.
\item Test exactness of $(2xy+3)\dd x+(x^2+4y)\dd y=0$ and find a potential.
\item Solve $y'=y^2$ with $y(0)=2$, including maximal interval.
\item Derive the analytic response of $C T'=-H(T-T_a)$ for constant $T_a$.
\end{enumerate}
\section{Higher order and transforms}
\begin{enumerate}[label=\textbf{P\arabic*.},start=7]
\item Solve $y''+4y=0$, $y(0)=0$, $y'(0)=2$.
\item Give the general solution of $y''-4y'+4y=0$.
\item Find a particular solution of $y''+y=\cos t$.
\item Solve $y'+y=1$ with $y(0)=0$ using Laplace transforms.
\item Find $\lap\{te^{-2t}\}(s)$ and its convergence half-plane.
\item Find the characteristic roots of $t^2y''+3ty'+y=0$ for $t>0$.
\end{enumerate}
\chapter{Original problem workbook: Part II}
\section{Systems, PDEs and boundary problems}
\begin{enumerate}[label=\textbf{P\arabic*.},start=13]
\item Solve $\mathbf x'=\operatorname{diag}(-1,-2)\mathbf x$ with $\mathbf x(0)=(1,3)^T$.
\item Find eigenvalues of $A=\begin{psmallmatrix}1&2\\2&1\end{psmallmatrix}$.
\item Find the trace/determinant classification for $A=\begin{psmallmatrix}-1&-1\\1&-1\end{psmallmatrix}$.
\item Solve the heat equation with fixed-zero endpoints and initial $\sin(\pi x/L)$.
\item Find all eigenvalues of $-y''=\lambda y$ on $0<x<1$, $y(0)=y(1)=0$.
\item Solve $-y''=2$ on $0<x<1$ with zero endpoints.
\item Give the compatibility condition for $y''+y=f(x)$ on $(0,\pi)$ with zero endpoints.
\item Derive $E'(t)$ for an undamped mass-spring model and interpret.
\end{enumerate}
\section{Existence, uniqueness and verification}
\begin{enumerate}[label=\textbf{P\arabic*.},start=21]
\item Determine why $y'=\sqrt{|y|}$, $y(0)=0$ need not have unique solutions.
\item Compute three Picard updates starting from $y_0(t)=1$ for $y'=y$, $y(0)=1$.
\item For $y'=\lambda y$, compare solutions with initial values $1$ and $1+\varepsilon$.
\item Approximate $y'=-y$ at $t=1$ with Euler $h=1/2$ and compute endpoint error.
\item Write a system and boundary residual vector for $y''=-2$, $y(0)=y(1)=0$.
\item For the standard Dirichlet Laplacian matrix on $N$ interior points, state eigenvalue scaling as $N\to\infty$ for fixed mode index.
\end{enumerate}
\chapter{Workbook solutions and verification cues}
\small
\begin{longtable}{p{0.10\textwidth}p{0.83\textwidth}}
\toprule Problem & Solution / verification criterion \\ \midrule\endhead
P1 & $y=3e^{t^2}$; $y'=2ty$. \\
P2 & $y=3(1-e^{-2t})$. \\
P3 & Equilibria $0,3$; positive derivative on $(0,3)$ and negative for $y>3$; $3$ attracts positive trajectories. \\
P4 & Exact: $M_y=2x=N_x$; $F=x^2y+3x+2y^2$. \\
P5 & $y=2/(1-2t)$, maximal interval $(-\infty,1/2)$. \\
P6 & $T=T_a+(T_0-T_a)e^{-Ht/C}$ for $C,H>0$. \\
P7 & $y(t)=\sin2t$. \\
P8 & $y=(C_1+C_2t)e^{2t}$. \\
P9 & $y_p=\tfrac12 t\sin t$. \\
P10 & $y=1-e^{-t}$. \\
P11 & $(s+2)^{-2}$ for $\operatorname{Re}(s)>-2$. \\
P12 & $(r+1)^2=0$; $y=t^{-1}(C_1+C_2\log t)$. \\
P13 & $\mathbf x=(e^{-t},3e^{-2t})^T$. \\
P14 & $3,-1$; check with eigenvectors $(1,1)$ and $(1,-1)$. \\
P15 & Trace $-2$, determinant $2$, discriminant $-4$: stable spiral. \\
P16 & $u(x,t)=e^{-\kappa\pi^2t/L^2}\sin(\pi x/L)$. \\
P17 & $\lambda_n=n^2\pi^2$, $n\ge1$. \\
P18 & $y=x(1-x)$. \\
P19 & $\int_0^\pi f(x)\sin x\dd x=0$. \\
P20 & For $E=\frac12mv^2+\frac12kx^2$, $E'=v(mv'+kx)=0$. \\
P21 & The right side lacks local Lipschitz continuity at $y=0$; zero and delayed-start solutions coexist. \\
P22 & $y_1=1+t$, $y_2=1+t+t^2/2$, $y_3=1+t+t^2/2+t^3/6$. \\
P23 & Difference $\varepsilon e^{\lambda t}$ for the same $\lambda$. \\
P24 & $y_2=(1-1/2)^2=1/4$; error $|1/4-e^{-1}|\approx0.11787944$. \\
P25 & $Y=(y,v)$, $Y'=(v,-2)$, $r=(Y_1(0),Y_1(1))$. \\
P26 & For fixed $n$, $\lambda_{n,h}\to n^2\pi^2$ as $h=1/(N+1)\to0$. \\
\bottomrule
\end{longtable}
\normalsize
\chapter{Laboratories, capstone integration, and study records}
\section{Lab A: analytic-versus-numeric convergence}
Benchmark $y'=-y$, $y(0)=1$ with Euler, RK4 and \texttt{solve\_ivp}. Use an identical $[0,1]$ interval and report endpoint errors; fit observed orders from a consistent hierarchy of steps. Provide source, package versions, plot, and assertions against the closed form.
\section{Lab B: two-tank conservation}
Set up $\mathbf x'=A\mathbf x$ with column sums of $A$ equal to zero and off-diagonal transfer rates nonnegative. Show conservation of $\mathbf 1^T\mathbf x$ and nonnegative concentrations. Investigate what changes when material is added, removed or reacted.
\section{Lab C: heat-equation mode analysis}
Implement the sine-series single mode and a finite-difference semi-discretization. Use energy decay and grid refinement as independent checks. Identify explicit-Euler stability restrictions and compare with an implicit solver.
\section{Lab D: Sturm--Liouville boundary spectrum}
Assemble the discrete Dirichlet Laplacian for at least three mesh sizes; compute low eigenvalues using a symmetric solver; compare against $(n\pi)^2$. Inspect normalized discrete eigenvectors and endpoint conformity.
\section{Lab E: stability certificate}
For a Hurwitz $A$, solve $A^TP+PA=-I$ for $P$ and verify positive definiteness. For a nonlinear system, state where a linear certificate is only local. Explain which parts of the stability claim are mathematical and which are numerical evidence.
\section{Final integrated capstone}
Build an original engineering case study involving a differential equation, at least one closed-form benchmark, an independent numerical model, parameter sensitivity, and verification under at least two operating scenarios. Deliver: governing derivation, assumptions register, dimensional audit, model-validity discussion, results table, numerical convergence study, limitations and an explicit distinction between textbook concepts and computational extensions.
\section{One-page study session template}
\begin{tabularx}{\textwidth}{>{\bfseries}p{0.26\textwidth}X}\toprule
Topic and chapter & \rule{0pt}{18pt}\\[8pt]
Prerequisites confirmed & \\[10pt]
Key equation and domain & \\[14pt]
Solution strategy & \\[14pt]
Proof or derivation & \\[14pt]
Analytic residual & \\[14pt]
Numerical comparison & \\[14pt]
Common failure mode & \\[14pt]
Next deliberate practice & \\[14pt]
\bottomrule\end{tabularx}

\appendix
\chapter{References and reproducibility}
\begin{enumerate}
\item R.~K. Nagle, E.~B. Saff, and A.~D. Snider, \emph{Fundamentals of Differential Equations and Boundary Value Problems}, 7th ed., Pearson, 2018, ISBN 978-0-321-97710-6. Primary source for the \emph{topic map and terminology}; cited locally by chapter and section.
\item SciPy project, \emph{solve\_ivp} reference, \url{https://docs.scipy.org/doc/scipy/reference/generated/scipy.integrate.solve_ivp.html} (modern extension).
\item SciPy project, \emph{solve\_bvp} reference, \url{https://docs.scipy.org/doc/scipy/reference/generated/scipy.integrate.solve_bvp.html} (modern extension).
\item SciPy project, \emph{eigsh} reference, \url{https://docs.scipy.org/doc/scipy/reference/generated/scipy.sparse.linalg.eigsh.html} (modern extension).
\item Python, NumPy, SciPy, and SymPy project documentation; computational checks require recording package versions and tolerances.
\end{enumerate}
\section*{Verification record template}
\begin{tabularx}{\textwidth}{>{\bfseries}p{0.30\textwidth}X}\toprule
Problem identifier & \rule{0pt}{15pt}\ \\ \midrule
ODE/PDE and domain & \\[8pt]
Initial/boundary conditions & \\[8pt]
Units and parameter values & \\[8pt]
Analytic residual & \\ [8pt]
Grid/step refinement & \\ [8pt]
Independent solver comparison & \\ [8pt]
Version and reproducibility notes & \\ [8pt]\bottomrule
\end{tabularx}
\end{document}
